% ============================================================================================
% BAB V IMPLEMENTASI
% MITIGASI HALUSINASI PADA RAG LLM MENGGUNAKAN
% QUERY REWRITING, CONTEXT RERANKING, DAN ACTIVE HALLUCINATION DETECTION
%
% TODO: Ganti semua placeholder [TODO: ...] dengan isi yang sebenarnya.
% TODO: Tambahkan screenshot atau listing kode yang relevan.
% ============================================================================================

\cleardoublepage
\chapter{IMPLEMENTASI}
\label{chap:implementasi}

Bab ini menjelaskan implementasi teknis dari rancangan sistem yang telah dijabarkan pada Bab~\ref{chap:perancangan}. Implementasi mencakup penyiapan lingkungan kerja, pembangunan indeks korpus PubMedQA, integrasi tiga komponen mitigasi halusinasi (ekspansi akronim, \textit{query rewriting}, dan \textit{context reranking}), serta penyiapan kerangka evaluasi. Seluruh komponen dikembangkan dalam bentuk \textit{notebook} Python yang dijalankan secara lokal pada satu unit \textit{laptop} pengembangan dan memanfaatkan beberapa layanan API eksternal untuk pemanggilan model bahasa berukuran besar.


% ============================================================================================
\section{Lingkungan Implementasi}
\label{sec:lingkungan-implementasi}
% ============================================================================================

Sistem mitigasi halusinasi RAG-LLM untuk sistem tanya jawab biomedis ini diimplementasikan pada satu mesin pengembangan tunggal dengan kombinasi infrastruktur lokal dan layanan eksternal. Infrastruktur lokal digunakan untuk menyimpan indeks vektor, menjalankan model \textit{embedding} dan \textit{reranker} ringan, serta mengeksekusi seluruh \textit{notebook} eksperimen. Layanan API eksternal dipakai untuk pemanggilan \textit{Large Language Model} (LLM) berukuran besar yang tidak dapat dijalankan secara lokal karena keterbatasan memori dan kapasitas komputasi. Subbab ini menjelaskan spesifikasi perangkat keras dan perangkat lunak, konfigurasi layanan API dan basis data, serta organisasi kode pada repositori penelitian.

\subsection{Spesifikasi Perangkat Keras dan Perangkat Lunak}
\label{subsec:spesifikasi-perangkat}

Pengembangan dan pengujian sistem dilakukan pada satu unit \textit{laptop} HP 14-ep1155TU sebagai mesin komputasi utama. Komputer ini menggunakan prosesor Intel Core 5 120U berkecepatan dasar 1{,}40 GHz dengan arsitektur x64, RAM 16 GB (15{,}7 GB \textit{usable}), serta \textit{Solid State Drive} (SSD) berkapasitas 477 GB. Mesin ini tidak memiliki akselerator \textit{Graphics Processing Unit} (GPU) terpisah; pengolahan grafis ditangani oleh kartu grafis terintegrasi Intel Graphics dengan memori \textit{shared} 128 MB. Oleh karena itu, seluruh beban komputasi lokal, termasuk pembangunan indeks BM25, pembentukan vektor \textit{embedding}, kueri ChromaDB, serta inferensi \textit{reranker}, dijalankan sepenuhnya di atas \textit{Central Processing Unit} (CPU). Pendekatan \textit{CPU-only} dipilih karena beban kerja sistem masih dapat diselesaikan dalam waktu yang wajar tanpa memerlukan GPU, mengingat ukuran korpus PubMedQA yang relatif kecil dan model \textit{reranker} yang dipilih sengaja berukuran ringan.

Dari sisi perangkat lunak, sistem berjalan di atas Windows 11 (64-bit). Bahasa pemrograman utama yang digunakan adalah Python versi 3.11 yang mendukung sebagian besar pustaka \textit{scientific computing} dan \textit{machine learning} terbaru, termasuk PyTorch, NumPy, pandas, dan turunannya. \textit{Notebook} eksperimen dijalankan menggunakan Jupyter Notebook melalui lingkungan virtual (\textit{virtualenv}) agar dependensi tiap modul tidak saling mengganggu. Sistem mengandalkan beberapa pustaka utama, yaitu \texttt{rank-bm25} untuk membangun indeks BM25, \texttt{chromadb} sebagai basis data vektor lokal, \texttt{sentence-transformers} untuk model \textit{reranker} berbasis \textit{cross-encoder}, serta klien resmi OpenAI, Anthropic, dan OpenRouter untuk pemanggilan LLM eksternal. Untuk evaluasi otomatis, sistem memakai pustaka \texttt{ragas} dengan versi kustom yang telah disesuaikan agar metrik \textit{faithfulness} dapat membaca konteks panjang tanpa pemotongan agresif. Model LLM lokal dijalankan menggunakan \texttt{Ollama} sebagai \textit{runtime} yang menyediakan antarmuka HTTP sederhana untuk model Llama versi terkuantisasi. Ringkasan seluruh komponen perangkat keras dan perangkat lunak disajikan pada Tabel~\ref{tbl:spesifikasi-implementasi}.

\begin{table}[H]
\centering
\caption{Ringkasan Perangkat Keras dan Perangkat Lunak yang Digunakan pada Penelitian}
\label{tbl:spesifikasi-implementasi}
\begin{tabular}{|p{5cm}|p{8.5cm}|}
\hline
\textbf{Komponen} & \textbf{Spesifikasi / Pustaka} \\ \hline
\multicolumn{2}{|l|}{\textit{Perangkat Keras}} \\ \hline
Mesin komputasi          & HP 14-ep1155TU \\ \hline
Prosesor                 & Intel Core 5 120U @ 1{,}40 GHz (x64) \\ \hline
RAM                      & 16 GB (15{,}7 GB \textit{usable}) \\ \hline
Penyimpanan              & SSD 477 GB \\ \hline
Kartu grafis             & Intel Graphics terintegrasi (128 MB \textit{shared}) \\ \hline
Mode eksekusi            & CPU-\textit{only} (tanpa akselerasi GPU) \\ \hline
\multicolumn{2}{|l|}{\textit{Perangkat Lunak}} \\ \hline
Sistem operasi           & Windows 11 (64-bit) \\ \hline
Bahasa pemrograman       & Python 3.11 \\ \hline
Lingkungan eksperimen    & Jupyter Notebook + \textit{virtualenv} \\ \hline
Pustaka BM25             & \texttt{rank-bm25} \\ \hline
Basis data vektor        & \texttt{chromadb} (mode lokal) \\ \hline
Pustaka \textit{embedding}   & \texttt{sentence-transformers}, OpenAI \texttt{text-embedding-3-small}, Ollama \texttt{nomic-embed-text} \\ \hline
Pustaka \textit{reranker}    & \texttt{sentence-transformers} dengan \texttt{ms-marco-MiniLM-L-6-v2} \\ \hline
\textit{Runtime} LLM lokal   & Ollama (\texttt{llama3.2}, \texttt{llama3.3}) \\ \hline
Klien LLM eksternal      & \texttt{openai}, \texttt{anthropic}, OpenRouter (kompatibel OpenAI) \\ \hline
Pustaka evaluasi         & \texttt{ragas} (kustom, dengan perbaikan truncation) \\ \hline
Pustaka pemrosesan data  & \texttt{pandas}, \texttt{numpy}, \texttt{datasets} (Hugging Face) \\ \hline
Pustaka NLP tambahan     & \texttt{nltk}, \texttt{regex}, \texttt{tiktoken} \\ \hline
\end{tabular}
\end{table}

\subsection{Konfigurasi Layanan API dan Basis Data}
\label{subsec:konfigurasi-api}

Sistem mengintegrasikan tiga layanan API model bahasa berbasis \textit{cloud}, yaitu OpenAI API untuk pemanggilan \texttt{gpt-4.1-mini} dan \texttt{text-embedding-3-small}, Anthropic API untuk pemanggilan \texttt{claude-haiku-4.5}, serta OpenRouter sebagai \textit{gateway} untuk mengakses \texttt{llama-3.3-70b-instruct} dari penyedia Meta. Pendekatan multi-penyedia ini dipilih agar setiap kategori model, yaitu \textit{closed-source} berkualitas tinggi (OpenAI, Anthropic) dan \textit{open-source} berukuran besar (Llama via OpenRouter), dapat dibandingkan secara adil pada beban kerja yang sama. Selain layanan eksternal, sistem juga memakai \texttt{Ollama} sebagai layanan lokal yang menjalankan model \texttt{llama3.2} dan \texttt{nomic-embed-text} secara \textit{on-premise} pada CPU, sehingga eksperimen tetap dapat berjalan tanpa biaya API ketika digunakan hanya untuk \textit{prototyping} atau pengujian awal.

Seluruh kredensial layanan eksternal, termasuk \texttt{OPENAI\_API\_KEY}, \texttt{ANTHROPIC\_API\_KEY}, dan \texttt{OPENROUTER\_API\_KEY}, disimpan pada berkas \texttt{.env} di akar repositori dan dimuat menggunakan pustaka \texttt{python-dotenv} pada awal setiap \textit{notebook}. Berkas \texttt{.env} ini dimasukkan ke dalam \texttt{.gitignore} agar informasi sensitif tidak ikut terunggah saat repositori didorong ke \textit{remote} GitHub. Untuk mencegah pelanggaran batas \textit{rate limit} pada masing-masing penyedia, setiap pemanggilan API diberi \textit{retry policy} sederhana dengan \textit{exponential backoff} pada kegagalan \textit{transient}, serta jeda kecil antarpermintaan ketika dijalankan pada eksperimen berukuran besar (500 pertanyaan). Parameter pemanggilan, seperti suhu (\textit{temperature}), \textit{top-p}, dan jumlah token maksimum, ditetapkan sama untuk seluruh konfigurasi agar perbandingan antarmodel bersifat \textit{apple-to-apple}.

Dari sisi basis data, sistem menggunakan dua jenis penyimpanan lokal yang saling melengkapi. Pertama, indeks BM25 dibangun menggunakan pustaka \texttt{rank-bm25} lalu disimpan dalam berkas \textit{pickle} (\texttt{pubmedqa\_bm25.pkl} dan \texttt{pubmedqa\_bm25\_sh.pkl}) sehingga indeks tidak perlu dibangun ulang setiap kali \textit{notebook} dijalankan. Kedua, indeks vektor menggunakan \texttt{ChromaDB} dalam mode \textit{persistent client} dengan direktori penyimpanan terpisah untuk masing-masing kombinasi praproses, yaitu \texttt{pubmedqa\_chroma/} (korpus asli), \texttt{pubmedqa\_chroma\_sh/} (korpus dengan ekspansi Schwartz-Hearst), serta \texttt{pubmedqa\_chroma\_llama/} (korpus dengan \textit{embedding} dari Ollama). Pemisahan direktori ini bertujuan memastikan setiap variasi eksperimen memiliki basis data yang konsisten dan dapat diuji ulang tanpa risiko \textit{cache poisoning} antarkonfigurasi. Ringkasan layanan API eksternal dan basis data lokal disajikan pada Tabel~\ref{tbl:layanan-api}.

\begin{table}[H]
\centering
\caption{Ringkasan Konfigurasi Layanan API Eksternal dan Basis Data Lokal}
\label{tbl:layanan-api}
\begin{tabular}{|p{4.5cm}|p{9cm}|}
\hline
\textbf{Layanan / Basis Data} & \textbf{Konfigurasi} \\ \hline
\multicolumn{2}{|l|}{\textit{Layanan API Model Bahasa Eksternal}} \\ \hline
OpenAI API           & Model \texttt{gpt-4.1-mini} untuk \textit{generation}, \texttt{text-embedding-3-small} (1.536d) untuk \textit{embedding}; kredensial \texttt{OPENAI\_API\_KEY} \\ \hline
Anthropic API        & Model \texttt{claude-haiku-4.5} untuk \textit{generation}; kredensial \texttt{ANTHROPIC\_API\_KEY} \\ \hline
OpenRouter           & Model \texttt{meta-llama/llama-3.3-70b-instruct} (kompatibel \textit{schema} OpenAI); kredensial \texttt{OPENROUTER\_API\_KEY} \\ \hline
\multicolumn{2}{|l|}{\textit{Layanan Lokal}} \\ \hline
Ollama               & \textit{Runtime} lokal untuk \texttt{llama3.2}, \texttt{llama3.3}, dan \texttt{nomic-embed-text} (768d); diakses melalui HTTP pada \texttt{localhost:11434} \\ \hline
\multicolumn{2}{|l|}{\textit{Basis Data Lokal}} \\ \hline
Indeks BM25          & Format \textit{pickle} (\texttt{pubmedqa\_bm25.pkl}, \texttt{pubmedqa\_bm25\_sh.pkl}) \\ \hline
ChromaDB             & Mode \textit{persistent client}, direktori terpisah per praproses (\texttt{pubmedqa\_chroma/}, \texttt{pubmedqa\_chroma\_sh/}, \texttt{pubmedqa\_chroma\_llama/}) \\ \hline
\multicolumn{2}{|l|}{\textit{Manajemen Konfigurasi}} \\ \hline
Berkas \texttt{.env} & Pusat penyimpanan kredensial API; dimuat melalui \texttt{python-dotenv}; dimasukkan ke \texttt{.gitignore} \\ \hline
\end{tabular}
\end{table}

\subsection{Struktur dan Organisasi Repositori}
\label{subsec:struktur-repo}

Kode sumber sistem disusun dengan prinsip pemisahan tanggung jawab agar tiap komponen utama dapat dikembangkan dan diuji secara independen. Struktur direktori repositori mencerminkan alur eksperimen, mulai dari penyiapan korpus, pembangunan indeks, pemanggilan komponen mitigasi, hingga visualisasi hasil pada antarmuka pengguna. Tabel~\ref{tbl:struktur-repo} merangkum direktori utama beserta fungsinya.

\begin{table}[H]
\centering
\caption{Struktur Direktori Utama Repositori Sistem}
\label{tbl:struktur-repo}
\begin{tabular}{|p{4.5cm}|p{9cm}|}
\hline
\textbf{Direktori / Berkas} & \textbf{Fungsi} \\ \hline
\texttt{notebooks/}                 & \textit{Notebook} eksperimen konfigurasi utama (\textit{baseline}, +QR, +CR, +QR+CR) untuk tiga model: Llama, OpenAI, dan Claude \\ \hline
\texttt{notebooks/BM25/}            & \textit{Notebook} variasi ablation BM25 saja (tanpa \textit{dense retrieval}) \\ \hline
\texttt{notebooks/BM25 Expansion/}  & \textit{Notebook} variasi dengan ekspansi akronim Schwartz-Hearst beserta skrip pembangunan indeks dan evaluasi \textit{retrieval} \\ \hline
\texttt{data/}                      & Dataset PubMedQA subset \textit{pqa\_labeled} hasil unduhan dari Hugging Face \\ \hline
\texttt{results/}                   & Hasil \textit{phase 1} (jawaban LLM dan label prediksi) dan \textit{phase 2} (skor RAGAS) per konfigurasi dalam format JSON, dengan subdirektori \texttt{BM25/} dan \texttt{BM25\_Expansion/} \\ \hline
\texttt{ui/}                        & Aplikasi web berbasis React + Vite + Tailwind untuk eksplorasi 500 pertanyaan hasil \textit{baseline} SH \\ \hline
\texttt{pyproject.toml}             & Definisi dependensi proyek dan konfigurasi \textit{virtualenv} \\ \hline
\texttt{.env}                       & Kredensial API eksternal (tidak diunggah ke repositori publik) \\ \hline
\end{tabular}
\end{table}

Pemisahan antara \textit{notebook} eksperimen utama, variasi ablation BM25 saja, dan variasi dengan ekspansi akronim, beserta direktori hasil per kelompok eksperimen, bertujuan agar setiap konfigurasi dapat dijalankan ulang secara mandiri tanpa risiko menimpa hasil eksperimen lainnya. Pendekatan ini juga memudahkan pelacakan saat melakukan analisis komparatif antarkonfigurasi pada Bab~\ref{chap:evaluasi}.


% ============================================================================================
\section{Implementasi Pipeline Pengolahan Data}
\label{sec:impl-pipeline-data}
% ============================================================================================

Tahap pengolahan data mengubah korpus mentah PubMedQA menjadi aset pengetahuan yang siap dipakai oleh sistem RAG, yaitu indeks BM25 dan indeks vektor pada ChromaDB. \textit{Pipeline} dirancang sebagai eksekusi sekali jalan yang dipanggil melalui \textit{notebook} pembangunan indeks, dan hasilnya disimpan ke berkas terstruktur sehingga tidak perlu diproses ulang setiap kali sistem dijalankan. Subbab ini menjelaskan tiga tahap \textit{pipeline}, yaitu pra-pemrosesan korpus, implementasi ekspansi akronim Schwartz-Hearst, serta pembangunan indeks BM25 dan vektor.

\subsection{Pra-pemrosesan Korpus PubMedQA}
\label{subsec:impl-praproses-korpus}

Korpus PubMedQA subset \textit{pqa\_labeled} dimuat menggunakan pustaka \texttt{datasets} dari Hugging Face dengan identitas dataset \texttt{qiaojin/PubMedQA} dan konfigurasi \texttt{pqa\_labeled}. Setiap entri dataset memiliki empat \textit{field} utama, yaitu \textit{question}, \textit{context}, \textit{long\_answer}, dan \textit{final\_decision}. Untuk keperluan pembangunan indeks, hanya \textit{field context} yang dipakai sebagai sumber dokumen; sedangkan \textit{question} dipakai sebagai \textit{query} pada saat \textit{retrieval}, dan dua \textit{field} terakhir disisihkan sebagai \textit{ground truth} evaluasi tanpa dimasukkan ke basis data agar tidak terjadi kebocoran jawaban (\textit{answer leakage}).

Fitur \textit{context} pada setiap entri berbentuk \textit{list} paragraf dengan label \textit{section} (\textit{BACKGROUND}, \textit{METHODS}, \textit{RESULTS}, \textit{CONCLUSIONS}, dan sebagainya). Pada implementasi, setiap paragraf diperlakukan sebagai satu \textit{chunk} terpisah dengan metadata berupa \textit{pubid}, \textit{section label}, dan posisi \textit{chunk} pada paper asal. Pendekatan ini dipilih karena ukuran paragraf pada abstrak biomedis sudah secara natural sesuai dengan \textit{context window} model \textit{embedding} dan tidak memerlukan strategi \textit{chunking} tambahan seperti \textit{sliding window} atau \textit{recursive splitting}. Setelah \textit{chunking} selesai, dilakukan normalisasi ringan berupa pembersihan \textit{whitespace} berlebih dan penghapusan karakter kontrol \textit{non-printable} jika ditemukan. Korpus dibangun dari seluruh 1.000 paper pada subset \textit{pqa\_labeled} dan menghasilkan kumpulan \textit{chunk} yang siap dipakai sebagai sumber dokumen.

\subsection{Implementasi Ekspansi Akronim Schwartz-Hearst}
\label{subsec:impl-ekspansi-sh}

Ekspansi akronim bertujuan menambahkan kepanjangan setiap akronim secara otomatis ke dalam teks dokumen, sehingga \textit{retriever} dapat menemukan dokumen yang relevan baik ketika pengguna mengetikkan akronim (\textit{Hyperbaric Oxygenation}) maupun bentuk panjangnya. Pada implementasi, langkah ini berjalan sekali pada saat penyusunan korpus dan hasilnya disimpan ke berkas indeks tersendiri.

Sistem bekerja dengan cara yang sederhana, yaitu memindai setiap kalimat pada dokumen untuk mencari pola tulisan ``\textit{kepanjangan (akronim)}'' yang umum dipakai pada paper biomedis. Sebagai contoh, ketika sistem menemukan frasa ``\textit{Hyperbaric Oxygenation} (HBO)'' pada salah satu paper, sistem akan mengenali bahwa HBO adalah akronim dari \textit{Hyperbaric Oxygenation} dan menyimpan pasangan tersebut ke kamus akronim. Pencocokan dilakukan dengan memeriksa apakah huruf-huruf akronim memang muncul di sepanjang kata pada bentuk panjangnya, sehingga akronim yang tidak konsisten (misalnya angka tahun atau singkatan kebetulan) tidak ikut tertangkap. Untuk akronim yang sering muncul namun tidak pernah didefinisikan langsung pada paper, sistem menyediakan kamus cadangan berisi 175 akronim biomedis yang paling umum sebagai \textit{fallback}.

Setelah pasangan akronim dan kepanjangannya teridentifikasi, sistem menjalankan substitusi pada seluruh \textit{chunk} dokumen pada paper yang sama. Setiap kemunculan akronim diubah menjadi format ``\texttt{AKRONIM (kepanjangan)}'' secara \textit{inline}, sehingga teks asli tetap terbaca dengan tambahan kepanjangan dalam tanda kurung. Pendekatan ini menjaga teks dokumen tidak banyak berubah, sekaligus memperkaya kosakata yang dipakai \textit{retriever}. Contoh hasil substitusi disajikan pada Tabel~\ref{tbl:ekspansi-contoh}.

\begin{table}[H]
\centering
\footnotesize
\setlength{\tabcolsep}{5pt}
\caption{Contoh Hasil Ekspansi Akronim pada Teks \textit{Chunk} Dokumen}
\label{tbl:ekspansi-contoh}
\begin{tabular}{p{1.3cm} p{5.2cm} p{6.2cm}}
\toprule
\textbf{Akronim} & \textbf{Teks Asli (Korpus Original)} & \textbf{Teks Setelah Ekspansi} \\
\midrule
HBO    & The mortality rate among the HBO-treated patients was 36\%. & The mortality rate among the HBO (hyperbaric oxygenation)-treated patients was 36\%. \\
\addlinespace
BMI    & Higher BMI was independently associated with disease severity. & Higher BMI (body mass index) was independently associated with disease severity. \\
\addlinespace
CHD    & The prevalence of CHD in this cohort was 14\%. & The prevalence of CHD (coronary heart disease) in this cohort was 14\%. \\
\addlinespace
APACHE & APACHE II score was used to assess illness severity at admission. & APACHE (Acute Physiology and Chronic Health Evaluation) II score was used to assess illness severity at admission. \\
\addlinespace
IBS    & Patients with IBS reported significantly more abdominal discomfort. & Patients with IBS (Irritable Bowel Syndrome) reported significantly more abdominal discomfort. \\
\bottomrule
\end{tabular}
\end{table}

Penerapan ekspansi pada korpus PubMedQA menambah panjang rata-rata \textit{chunk} dari 423 karakter menjadi 473 karakter, namun jumlah total \textit{chunk} pada korpus tidak berubah karena ekspansi bersifat substitusi \textit{inline}. Inti dari proses pencocokan huruf akronim dengan kandidat kepanjangan dijabarkan pada Algoritma~\ref{alg:sh-validation}.

\begin{algorithm}[H]
\caption{Validasi Pasangan Akronim-Kepanjangan dengan Schwartz-Hearst}
\label{alg:sh-validation}
\begin{algorithmic}[1]
\Require Akronim $s$ dengan panjang $n$, kandidat kepanjangan $l$ dengan panjang $m$
\Ensure Pasangan kepanjangan valid atau $\bot$ apabila tidak cocok
\State $i \gets n - 1$ \Comment{indeks huruf terakhir akronim}
\State $j \gets m - 1$ \Comment{indeks karakter terakhir kandidat}
\While{$i \geq 0$}
  \If{$i = 0$} \Comment{huruf pertama akronim wajib di awal kata}
    \While{$j \geq 0$ \textbf{and not} ($l[j] = s[i]$ \textbf{and} \Call{awalKata}{$l, j$})}
      \State $j \gets j - 1$
    \EndWhile
  \Else \Comment{huruf lain boleh dari posisi internal kata}
    \While{$j \geq 0$ \textbf{and} $l[j] \neq s[i]$}
      \State $j \gets j - 1$
    \EndWhile
  \EndIf
  \If{$j < 0$} \Return $\bot$ \EndIf
  \State $i \gets i - 1$
  \State $j \gets j - 1$
\EndWhile
\State \Return $l[j+1 \ldots m-1]$ \Comment{substring kepanjangan tervalidasi}
\end{algorithmic}
\end{algorithm}

\subsection{Pembangunan Indeks BM25 dan Vektor}
\label{subsec:impl-build-index}

Setelah \textit{chunk} dokumen disiapkan, sistem membangun dua indeks pencarian yang saling melengkapi, yaitu indeks BM25 untuk pencarian berbasis kata kunci dan indeks vektor untuk pencarian berbasis kemiripan makna. Kedua indeks dibangun sekali pada awal eksperimen dan disimpan ke berkas terpisah, sehingga eksekusi berikutnya cukup memuat indeks yang sudah jadi tanpa perlu membangun ulang dari awal.

Untuk indeks BM25, setiap \textit{chunk} terlebih dahulu ditokenisasi dengan cara yang sederhana, yaitu memisahkan teks berdasarkan tanda baca dan spasi, lalu mengubah seluruh karakter menjadi huruf kecil. Hasil tokenisasi seluruh \textit{chunk} membentuk korpus terindeks yang siap menerima kueri pengguna pada tahap pencarian. Pada implementasi, sistem menyimpan dua versi indeks BM25 secara terpisah, yaitu versi yang dibangun dari korpus asli dan versi yang dibangun dari korpus dengan ekspansi akronim, agar eksperimen perbandingan dapat dijalankan tanpa mencampur kedua kondisi.

Indeks vektor dibangun menggunakan basis data vektor ChromaDB. Setiap \textit{chunk} diubah menjadi representasi vektor numerik melalui pemanggilan model \textit{embedding}, lalu disimpan beserta metadata berupa identitas paper sumber, label \textit{section}, dan posisi \textit{chunk}. Pasangan model \textit{embedding} disesuaikan dengan LLM yang dipakai pada tahap \textit{generation}, yaitu \texttt{text-embedding-3-small} dari OpenAI (1.536 dimensi) untuk pasangan dengan GPT-4.1-mini dan Claude Haiku 4.5, dan \texttt{nomic-embed-text} yang dijalankan secara lokal melalui Ollama (768 dimensi) untuk pasangan dengan Llama. Pemisahan model \textit{embedding} ini menjaga konsistensi internal antar konfigurasi sekaligus memastikan eksperimen dengan Llama tidak bergantung pada layanan API eksternal.


% ============================================================================================
\section{Implementasi Komponen Retrieval}
\label{sec:impl-retrieval}
% ============================================================================================

Komponen \textit{retrieval} bertugas mengambil \textit{chunk} yang paling relevan terhadap pertanyaan pengguna dari kedua basis data yang telah dibangun pada tahap sebelumnya. Implementasi menggunakan pendekatan \textit{hybrid retrieval} yang menggabungkan BM25 dan \textit{dense retrieval} melalui metode \textit{Reciprocal Rank Fusion} (RRF). Subbab ini menjelaskan implementasi ketiga komponen tersebut.

\subsection{BM25}
\label{subsec:impl-bm25}

Komponen BM25 bertugas mencari \textit{chunk} dokumen yang paling banyak memuat istilah kueri, sebagaimana telah dijelaskan secara teoritis pada Bab~\ref{chap:studi-literatur}. Pada implementasi, sistem memuat indeks BM25 yang sudah dibangun pada tahap sebelumnya, lalu menjalankan kueri pengguna melalui dua langkah, yaitu menyamakan format kueri dengan format korpus melalui tokenisasi, dan menghitung skor BM25 antara kueri tersebut dengan seluruh \textit{chunk} pada indeks.

Tokenisasi pada kueri dilakukan dengan prosedur yang sama dengan tokenisasi korpus, yaitu memisahkan teks berdasarkan tanda baca dan spasi serta mengubah seluruh karakter menjadi huruf kecil. Kesamaan prosedur ini penting agar kueri seperti \textit{``Does HBO improve survival?''} terurai menjadi token \textit{[does, hbo, improve, survival]} yang sama persis dengan token pada korpus, sehingga pencocokan kata kunci tidak terganggu oleh perbedaan kapitalisasi atau tanda baca. Setelah tokenisasi, sistem menghitung skor BM25 antara token kueri dan setiap \textit{chunk} menggunakan parameter standar pada literatur, yaitu $k_1 = 1{,}5$ untuk pengaturan saturasi \textit{term frequency} dan $b = 0{,}75$ untuk pengaturan normalisasi panjang dokumen. Setiap \textit{chunk} memperoleh satu skor agregat, lalu seluruh \textit{chunk} diurutkan dari skor tertinggi ke terendah. Lima puluh \textit{chunk} teratas (\textit{BM25 Top-50}) kemudian dipilih sebagai kandidat untuk tahap fusi peringkat dengan hasil \textit{dense retrieval}.

Untuk memberikan gambaran intuitif tentang cara kerja BM25, Tabel~\ref{tbl:bm25-contoh} memperlihatkan ilustrasi pencocokan kueri \textit{``hyperbaric oxygenation mortality treatment''} (yang diturunkan dari paper PubMedQA tentang efek terapi oksigen hiperbarik) terhadap tiga \textit{chunk} kandidat dengan karakteristik berbeda. Skor pada tabel ini dihitung secara sederhana sebagai ilustrasi untuk menjelaskan kontribusi setiap komponen BM25 dan bukan merupakan skor sebenarnya dari korpus penelitian.

\begin{table}[H]
\centering
\footnotesize
\setlength{\tabcolsep}{4pt}
\caption{Ilustrasi Skor BM25 untuk Kueri ``hyperbaric oxygenation mortality treatment'' pada Tiga \textit{Chunk} Kandidat}
\label{tbl:bm25-contoh}
\begin{tabular}{p{0.7cm} p{5.6cm} cccc}
\toprule
\textbf{ID} & \textbf{Cuplikan \textit{Chunk}} & \textbf{Panjang} & \textbf{Match} & \textbf{IDF} & \textbf{Skor} \\
\midrule
$C_1$ & The mortality rate among the HBO (hyperbaric oxygenation)-treated patients was 36\%, lower than the control group. & 18 & 4 & tinggi & \textbf{8{,}21} \\
\addlinespace
$C_2$ & Patients underwent hyperbaric oxygenation treatment for chronic wound healing in a prospective study. & 13 & 3 & tinggi & 6{,}54 \\
\addlinespace
$C_3$ & The treatment was effective in reducing patient mortality across various clinical conditions. & 14 & 2 & rendah & 2{,}11 \\
\bottomrule
\end{tabular}
\end{table}

Pada Tabel~\ref{tbl:bm25-contoh}, kolom \textit{Match} menunjukkan jumlah istilah kueri yang muncul pada \textit{chunk}, kolom \textit{IDF} menunjukkan kekayaan istilah langka (\textit{hyperbaric, oxygenation, HBO}) dibandingkan istilah umum (\textit{treatment, mortality}), dan kolom \textit{Skor} adalah agregasi sesuai Persamaan~\ref{eq:bm25}. \textit{Chunk} $C_1$ mendapat skor tertinggi karena memuat hampir seluruh istilah kueri, termasuk dua istilah langka bernilai IDF tinggi (\textit{hyperbaric, oxygenation, mortality}). \textit{Chunk} $C_2$ mendapat skor menengah karena memuat istilah langka yang sama namun melewatkan istilah \textit{mortality}. \textit{Chunk} $C_3$ mendapat skor paling rendah meskipun memuat \textit{treatment} dan \textit{mortality}, karena kedua istilah tersebut umum di seluruh korpus sehingga bobot IDF-nya kecil. Contoh ini juga memperlihatkan mengapa ekspansi akronim memberi pengaruh signifikan pada BM25; tanpa ekspansi, \textit{chunk} yang menulis \textit{HBO} saja tanpa kepanjangan tidak akan tertangkap oleh kueri yang menyebut \textit{hyperbaric oxygenation}, sehingga peluang dokumen relevan ditemukan menjadi berkurang.


\subsection{Dense Retrieval}
\label{subsec:impl-dense}

Komponen \textit{dense retrieval} mencari \textit{chunk} dokumen berdasarkan kemiripan makna dengan kueri pengguna, sebagaimana telah dijelaskan secara teoritis pada Bab~\ref{chap:studi-literatur}. Pada implementasi, sistem mengubah kueri pengguna menjadi vektor melalui model \textit{embedding}, lalu meminta basis data vektor untuk mengembalikan \textit{chunk} dengan vektor yang paling dekat dengan vektor kueri tersebut.

Pemilihan model \textit{embedding} disesuaikan dengan LLM yang dipakai pada tahap \textit{generation}. Untuk konfigurasi yang memakai GPT-4.1-mini atau Claude Haiku 4.5, sistem memanggil model \texttt{text-embedding-3-small} dari OpenAI yang menghasilkan vektor berdimensi 1.536; sedangkan untuk konfigurasi yang memakai Llama, sistem memanggil model \texttt{nomic-embed-text} yang dijalankan secara lokal melalui Ollama dan menghasilkan vektor berdimensi 768. Penggunaan model \textit{embedding} yang sama saat \textit{indexing} dan saat kueri adalah syarat wajib, karena dua model yang berbeda akan menempatkan teks pada ruang vektor yang berbeda pula sehingga perhitungan kemiripan tidak bermakna. Setelah vektor kueri terbentuk, basis data ChromaDB menghitung \textit{cosine similarity} antara vektor tersebut dengan seluruh vektor pada koleksi, lalu mengembalikan 50 \textit{chunk} teratas (\textit{Dense Top-50}) sebagai kandidat untuk tahap fusi peringkat.

Untuk memberikan gambaran tentang kekuatan \textit{dense retrieval} dalam menangkap kemiripan makna, Tabel~\ref{tbl:dense-contoh} memperlihatkan ilustrasi pencocokan kueri \textit{``Does physical activity reduce postnatal depression?''} (yang merupakan parafrasa dari pertanyaan PubMedQA mengenai olahraga saat hamil dan depresi pasca-melahirkan) terhadap tiga \textit{chunk} kandidat. Nilai \textit{cosine similarity} pada tabel ini dihitung secara sederhana sebagai ilustrasi dan bukan merupakan skor sebenarnya dari korpus penelitian.

\begin{table}[H]
\centering
\footnotesize
\setlength{\tabcolsep}{4pt}
\caption{Ilustrasi Skor \textit{Cosine Similarity} untuk Kueri ``Does physical activity reduce postnatal depression?'' pada Tiga \textit{Chunk} Kandidat}
\label{tbl:dense-contoh}
\begin{tabular}{p{0.7cm} p{8.2cm} p{1.8cm} c}
\toprule
\textbf{ID} & \textbf{Cuplikan \textit{Chunk}} & \textbf{Kemiripan Makna} & \textbf{Skor} \\
\midrule
$C_1$ & Regular exercise during pregnancy was associated with lower depression scores after delivery. & sangat dekat (parafrasa langsung) & \textbf{0{,}872} \\
\addlinespace
$C_2$ & Aerobic training showed mood improvement in postpartum women across the trial cohort. & dekat (sinonim multi-istilah) & 0{,}791 \\
\addlinespace
$C_3$ & Pregnancy outcomes were measured using standard fetal monitoring techniques in the study. & jauh (hanya istilah ``pregnancy'') & 0{,}312 \\
\bottomrule
\end{tabular}
\end{table}

Pada Tabel~\ref{tbl:dense-contoh}, \textit{chunk} $C_1$ memperoleh skor tertinggi karena memuat parafrasa langsung dari setiap konsep pada kueri, yaitu \textit{exercise} untuk \textit{physical activity}, \textit{depression} untuk \textit{depression}, dan \textit{after delivery} untuk \textit{postnatal}. \textit{Chunk} $C_2$ memperoleh skor menengah karena memakai sinonim yang lebih jauh (\textit{aerobic training} untuk \textit{physical activity}, \textit{mood improvement} untuk \textit{reduce depression}, \textit{postpartum} untuk \textit{postnatal}), namun maknanya tetap konsisten dengan kueri sehingga vektornya tetap dekat dengan vektor kueri. \textit{Chunk} $C_3$ memperoleh skor rendah karena meskipun memuat istilah \textit{pregnancy}, isinya membahas pemantauan janin sehingga maknanya jauh berbeda dengan kueri tentang depresi.

Ilustrasi ini menjelaskan mengapa \textit{dense retrieval} bersifat komplementer terhadap BM25. Pada kueri yang memuat banyak parafrasa dan sinonim seperti contoh di atas, BM25 hanya mampu menangkap \textit{chunk} $C_3$ karena kemunculan istilah \textit{pregnancy} secara eksak, padahal $C_3$ justru kurang relevan. \textit{Dense retrieval} berhasil mengangkat $C_1$ dan $C_2$ yang sebenarnya jauh lebih relevan, walaupun kosakatanya berbeda dengan kueri. Kombinasi kedua jenis \textit{retrieval} pada tahap \textit{hybrid retrieval} kemudian menjaga agar dokumen yang ditemukan mencakup kedua kategori, yaitu dokumen yang cocok berdasarkan kata kunci spesifik maupun dokumen yang cocok berdasarkan kemiripan makna.

\subsection{Hybrid Retrieval dengan Reciprocal Rank Fusion}
\label{subsec:impl-hybrid-rrf}

Hasil dari kedua \textit{retriever} digabungkan melalui fungsi \texttt{rrf\_merge} yang mengimplementasikan rumus \textit{Reciprocal Rank Fusion} sesuai \parencite{cormack2009reciprocal}. Untuk setiap \textit{chunk} yang muncul pada salah satu atau kedua daftar hasil, sistem menghitung skor RRF sebagai $\sum_{i} \frac{1}{k + \text{rank}_i}$ dengan parameter $k = 60$ sesuai rekomendasi penulis asli, dan $\text{rank}_i$ adalah posisi \textit{chunk} pada daftar ke-$i$. \textit{Chunk} yang muncul pada kedua daftar akan mendapat kontribusi skor dari keduanya, sehingga dokumen yang ditemukan oleh BM25 maupun \textit{dense retrieval} secara konsisten akan terangkat ke peringkat atas. Setelah skor RRF dihitung, 20 \textit{chunk} dengan skor tertinggi (\textit{Top-20}) dipilih sebagai kandidat untuk tahap berikutnya, yaitu \textit{Context Reranking} atau langsung tahap \textit{generation}. Prosedur lengkap penggabungan kedua daftar peringkat dijabarkan pada Algoritma~\ref{alg:rrf}.

\begin{algorithm}[H]
\caption{Penggabungan Peringkat dengan \textit{Reciprocal Rank Fusion} (RRF)}
\label{alg:rrf}
\begin{algorithmic}[1]
\Require Daftar peringkat BM25 $R_B$ dan \textit{dense} $R_D$; konstanta $k = 60$; jumlah keluaran $N = 20$
\Ensure Daftar Top-$N$ \textit{chunk} berdasarkan skor RRF
\State $S \gets \{\}$ \Comment{peta dari \textit{chunk\_id} ke skor RRF}
\ForAll{$R_i \in \{R_B, R_D\}$}
  \For{$r \gets 1$ \textbf{to} $|R_i|$}
    \State $c \gets R_i[r]$ \Comment{\textit{chunk} pada peringkat ke-$r$}
    \State $S[c] \gets S[c] + \dfrac{1}{k + r}$ \Comment{akumulasi kontribusi tiap \textit{retriever}}
  \EndFor
\EndFor
\State $L \gets$ urutkan kunci $S$ berdasarkan nilai secara menurun
\State \Return $L[1 \ldots N]$
\end{algorithmic}
\end{algorithm}


% ============================================================================================
\section{Implementasi Komponen Mitigasi Halusinasi}
\label{sec:impl-mitigasi}
% ============================================================================================

Pada subbab ini dijelaskan implementasi dua komponen mitigasi yang bekerja di sekitar tahap \textit{retrieval}, yaitu \textit{Query Rewriting} (QR) yang bekerja sebelum \textit{retrieval} dan \textit{Context Reranking} (CR) yang bekerja sesudahnya. Komponen ketiga, yaitu ekspansi akronim Schwartz-Hearst, telah dijelaskan pada Subbab~\ref{subsec:impl-ekspansi-sh} karena sifatnya berupa praproses dokumen, bukan modifikasi alur \textit{retrieval}.

\subsection{Implementasi Query Rewriting}
\label{subsec:impl-qr}

Komponen \textit{Query Rewriting} diimplementasikan sebagai fungsi yang menerima pertanyaan asli pengguna dan mengembalikan versi yang telah ditulis ulang menjadi lebih kaya akan istilah medis. Fungsi ini mengirimkan \textit{prompt} khusus kepada LLM yang berisi instruksi untuk menghasilkan ulang pertanyaan dengan tiga ketentuan, yaitu (1) memperluas istilah medis dengan menyertakan sinonim atau bentuk akronim, (2) mengeksplisitkan entitas biomedis yang implisit pada pertanyaan asli, dan (3) mempertahankan makna serta jenis pertanyaan agar jawaban yang dihasilkan tetap relevan terhadap kebutuhan pengguna.

LLM yang dipakai untuk \textit{Query Rewriting} adalah model yang sama dengan model untuk tahap \textit{generation} pada konfigurasi bersangkutan (\textit{gpt-4.1-mini}, \textit{claude-haiku-4.5}, atau \textit{llama-3.3-70b-instruct}). Pendekatan ini menjaga kompatibilitas gaya bahasa antara penulisan ulang dan pembentukan jawaban. Hasil \textit{rewriting} menggantikan pertanyaan asli pada tahap \textit{retrieval} (baik BM25 maupun \textit{dense}), namun pertanyaan asli tetap diteruskan ke tahap \textit{generation} agar LLM menjawab pertanyaan yang sebenarnya diajukan pengguna, bukan versi tulis-ulang internal.

\subsection{Implementasi Context Reranking}
\label{subsec:impl-cr}

Komponen \textit{Context Reranking} diimplementasikan menggunakan model \textit{cross-encoder} \texttt{ms-marco-MiniLM-L-6-v2} dari pustaka \texttt{sentence-transformers}. Model dipilih karena ukurannya kecil (sekitar 22 MB), dapat berjalan cepat pada CPU tanpa GPU, dan terbukti efektif untuk tugas \textit{passage reranking}. Berbeda dengan model \textit{embedding} yang menghasilkan satu vektor per teks, \textit{cross-encoder} menerima pasangan (\textit{query}, \textit{passage}) sekaligus pada satu inferensi dan menghasilkan skor relevansi tunggal, sehingga kualitas penilaian lebih akurat untuk tugas \textit{re-ranking} skala kecil.

Pada implementasi, 20 \textit{chunk} teratas hasil RRF (\textit{Top-20}) dipasangkan satu per satu dengan kueri pengguna, sehingga terbentuk 20 pasangan masukan. Pasangan tersebut diteruskan ke metode \texttt{predict} dari objek \texttt{CrossEncoder} secara \textit{batch} agar inferensi efisien. Setelah skor relevansi diperoleh, \textit{chunk} diurutkan ulang dari skor tertinggi ke terendah, dan lima \textit{chunk} paling relevan (\textit{Top-5}) dipilih sebagai konteks akhir yang diteruskan ke tahap \textit{generation}. Pemilihan \textit{Top-5} ini konsisten dengan rancangan pada Bab~\ref{chap:perancangan} dan menjadi standar input konteks untuk seluruh konfigurasi sistem.


% ============================================================================================
\section{Implementasi Komponen Generation dan Kerangka Evaluasi}
\label{sec:impl-generation-evaluasi}
% ============================================================================================

\subsection{Implementasi Tahap Generation}
\label{subsec:impl-generation}

Tahap \textit{generation} menerima pertanyaan asli pengguna beserta lima \textit{chunk} konteks pilihan (\textit{Top-5}), lalu meminta LLM membentuk jawaban berupa label klasifikasi (\textit{yes}, \textit{no}, atau \textit{maybe}) beserta penjelasan singkat yang berlandaskan pada konteks. \textit{Prompt} yang dikirim ke LLM disusun dengan tiga komponen, yaitu (1) instruksi peran sebagai asisten peneliti biomedis yang menjawab berdasarkan dokumen yang diberikan, (2) lima \textit{chunk} konteks yang dipisahkan oleh delimiter \texttt{===} dan dilengkapi dengan penanda sumber paper, serta (3) pertanyaan pengguna disertai instruksi eksplisit agar jawaban diakhiri dengan satu dari tiga label pada baris terakhir.

Pemanggilan LLM dilakukan melalui klien resmi masing-masing penyedia, yaitu \texttt{openai} untuk OpenAI dan OpenRouter (\textit{schema}-nya kompatibel), \texttt{anthropic} untuk Claude, serta antarmuka HTTP Ollama untuk Llama lokal. Parameter pemanggilan ditetapkan sama untuk seluruh model, yaitu \textit{temperature} sebesar 0 untuk menjamin determinisme jawaban dan jumlah token maksimum yang disesuaikan dengan kebutuhan jawaban biomedis singkat. Keluaran LLM berupa teks bebas kemudian diproses oleh fungsi ekstraksi label yang mencoba tiga strategi secara berurutan, yaitu (1) pencocokan langsung pada baris terakhir, (2) penelusuran kata kunci \textit{yes}, \textit{no}, atau \textit{maybe} pada seluruh teks, dan (3) pencocokan berbasis ekspresi reguler sebagai penanganan kasus tepi. Apabila ketiga strategi gagal, label ditetapkan sebagai \textit{unknown} dan ditandai sebagai sampel yang tidak dapat dievaluasi.

\subsection{Implementasi Kerangka Evaluasi}
\label{subsec:impl-evaluasi}

Evaluasi sistem dijalankan dalam dua fase yang menyimpan hasilnya secara inkremental untuk menghindari kehilangan data apabila proses terhenti di tengah jalan. Fase pertama (\textit{phase 1}) menjalankan seluruh proses \textit{retrieval} dan \textit{generation}, kemudian menyimpan jawaban LLM beserta label prediksi dan lima \textit{chunk} konteks ke berkas \texttt{*\_phase1\_answers.json}. Fase kedua (\textit{phase 2}) memuat keluaran fase pertama dan menjalankan evaluasi RAGAS pada empat metrik, yaitu \textit{faithfulness}, \textit{context recall}, \textit{answer relevancy}, dan \textit{context precision}, lalu menyimpannya ke berkas \texttt{*\_phase2\_custom.json}. Pemisahan fase ini juga memungkinkan eksperimen yang sama dijalankan ulang dengan metrik berbeda tanpa perlu memanggil ulang LLM untuk membentuk jawaban.

Implementasi evaluasi RAGAS pada penelitian ini menjalankan empat metrik tersebut pada setiap pasangan (\textit{question}, \textit{contexts}, \textit{answer}, \textit{ground truth}) yang dihasilkan pada fase pertama. Setiap konteks yang diteruskan ke evaluator merupakan lima \textit{chunk} hasil \textit{retrieval} yang sama persis dengan yang dipakai LLM saat membentuk jawaban, sehingga skor evaluasi mencerminkan kualitas jawaban berdasarkan informasi yang benar-benar dilihat oleh sistem.

Selain evaluasi RAGAS, sistem juga mengukur kualitas \textit{retrieval} secara terpisah menggunakan tiga metrik \textit{Information Retrieval} (IR) standar, yaitu \textit{Recall@5}, \textit{MAP@5}, dan \textit{nDCG@5}. Implementasi metrik ini berada pada modul \texttt{evaluate\_retrieval.py} pada direktori \texttt{notebooks/BM25 Expansion/}. Modul ini membandingkan lima \textit{chunk} teratas hasil \textit{retrieval} dengan \textit{chunk} yang memuat \textit{long\_answer} sebagai \textit{ground truth} relevansi, lalu menghitung ketiga metrik untuk seluruh sampel uji. Hasil pengukuran disimpan ke berkas \texttt{retrieval\_metrics.json} dan dipakai untuk membandingkan kualitas \textit{retrieval} antara konfigurasi BM25 saja, BM25 + \textit{dense} (Hybrid), serta variasi dengan dan tanpa ekspansi akronim.
